% !TEX root = ../main-es.tex
% Las cajas de ayuda no son contenido: texcount las ignora (ver `make wordcount`).
%TC:macro \limite [ignore,ignore]
%TC:macro \nota [ignore]
%TC:macro \TODO [ignore]
%TC:macro \parencite [ignore]
%TC:macro \textcite [ignore]
\section{Aportes de las referencias principales}
\limite{0}{Aportes de las referencias principales: seleccione dos referencias clave y
describa en qué medida contribuyen a la construcción del proyecto (por ejemplo:
fundamentos teóricos, antecedentes, metodologías, aplicaciones).}

\subsection*{\textcite{russellNorvig2020}: Fundamentos teóricos}

Es el texto de la asignatura y a la vez el andamiaje conceptual completo del proyecto,
pues de él salen las tres piezas que lo estructuran. Primero, la taxonomía de agentes y
entornos del capítulo 2 fija las seis arquitecturas comparadas así como la clasificación
del entorno de tarea (parcialmente observable, estocástico, secuencial y dinámico).
Luego, la búsqueda informada del capítulo 3 aporta el A* que el agente invoca en cada
reescritura (contenido del curso usado como motor funcional y no solo como cita).
Finalmente, la planeación en línea del capítulo 11 (monitoreo de acción, de plan y de
objetivo) define qué significa detectar que un plan dejó de ser válido, que es
precisamente el evento que dispara toda la dinámica de replanificación.

\subsection*{\textcite{gao2026selfevolving}: La agenda del campo}

Si R\&N aporta el andamiaje conceptual, el survey de agentes auto-evolutivos aporta la
pertinencia y el interlocutor, pues es el trabajo que le pone nombre y agenda a la clase
de sistemas que el proyecto estudia. Su contribución es doble. Primero, organiza el campo
en cuatro preguntas (\emph{what}, \emph{when}, \emph{how} y \emph{where to evolve}), de
forma que la pregunta de este proyecto no queda como una curiosidad de ingeniería sino
como la dimensión \emph{when} de una agenda explícita, que el propio survey reconoce
abierta. Luego, y esto es lo que define el aporte metodológico, taxonomiza los mecanismos
de auto-modificación sin entregar criterio de decisión: Describe con qué se evoluciona y
cómo, pero no bajo qué condiciones conviene hacerlo, por lo que la frontera medible en el
plano $(\lambda, c)$ es precisamente la operacionalización que falta.

De ese encuadre se desprenden además las instancias que sostienen el Cuadro~1, dado que el
survey reúne en una sola clase sistemas que de otro modo se leerían como anécdotas
inconexas (bibliotecas de habilidades programadas por el propio agente, búsqueda de flujos
representados en código, reescritura de memoria tras el fracaso), lo cual permite defender
que las seis arquitecturas comparadas son abstracciones controladas de sistemas reales y
no construcciones ad hoc.

Cabe anotar de dónde sale el vocabulario de costo con el que se formula la hipótesis, pues
no lo aporta el survey sino la tradición de metarazonamiento y racionalidad acotada
\parencite{zilberstein2011metareasoning}: Deliberar es una acción con costo que compite con
actuar, de modo que el costo $c$ de edición y el precio $\kappa$ por tick de deliberación
operacionalizan ese costo, el presupuesto $B$ del agente gobernado es una política de
deliberación explícita y el \emph{thrashing} medido por las ediciones desperdiciadas es el
modo de fallo clásico de una política mal calibrada.
